\documentclass[hidelinks]{article}
\usepackage[a4paper, total={7in, 10in}]{geometry}
\usepackage[dvipsnames]{xcolor}
\usepackage{amsmath, amssymb, amsthm}
\usepackage{tikz}
\usepackage{tkz-euclide}
\usepackage[unicode]{hyperref}
\usepackage[all]{hypcap}
\usepackage{fancyhdr}
\usepackage{listings}
\usepackage{subcaption}
\usepackage{enumitem}
\usepackage{graphicx}

\lstset{basicstyle=\ttfamily\small,columns=fullflexible,frame=single,breaklines=true,showstringspaces=false}

\title{\textbf{Dielectrics in Electrostatics Test}}
\author{Diego Juarez}
\date{March 25th, 2026}

\pagestyle{fancy}
\fancyhf{}
\fancyhead[L]{Dielectrics Test}
\fancyhead[R]{Jackson-style Practice}
\fancyfoot[C]{\thepage}

\begin{document}
\hypersetup{bookmarksnumbered=true}
\pagecolor{black}
\color{white}
\maketitle

\begin{Large}
\tableofcontents
\end{Large}
\pagebreak

\section{Source Notebook}
This problem set follows the notebook
\url{../notebooks/Dielectrics in electrostatics.ipynb}
and focuses on polarization, bound charge, displacement field, interfaces, and dielectric spheres.

\section{Basic Definitions}
\subsection{Problem 1: Polarization and Bound Charge}
Define the polarization field $\mathbf{P}(\mathbf{r})$. Show that the bound volume charge density is
\[
\rho_b=-\nabla\cdot \mathbf{P}.
\]
State also the expression for bound surface charge density.

\subsection{Problem 2: Displacement Field}
Define the electric displacement field
\[
\mathbf{D}=\varepsilon_0\mathbf{E}+\mathbf{P}.
\]
Starting from Gauss's law, derive
\[
\nabla\cdot \mathbf{D}=\rho_f,
\]
where $\rho_f$ is the free charge density.

\subsection{Problem 3: Linear Dielectrics}
For a linear isotropic dielectric, write the constitutive relation between $\mathbf{P}$ and $\mathbf{E}$. Then derive $\mathbf{D}=\varepsilon \mathbf{E}$ and identify $\varepsilon$ in terms of susceptibility.

\section{Canonical Configurations}
\subsection{Problem 4: Uniformly Polarized Slab}
Consider a slab polarized uniformly in the $+\hat{\mathbf{z}}$ direction.
\begin{enumerate}[label=\alph*)]
\item Find the bound volume charge density.
\item Find the bound surface charge densities on the two faces.
\item Sketch or describe the qualitative electric field produced by the slab.
\end{enumerate}

\subsection{Problem 5: Uniformly Polarized Sphere}
A sphere of radius $R$ has uniform polarization $\mathbf{P}=P_0\hat{\mathbf{z}}$.
\begin{enumerate}[label=\alph*)]
\item Find $\rho_b$ and $\sigma_b$.
\item Explain why the interior electric field is uniform.
\item State the direction of the interior field relative to $\mathbf{P}$.
\end{enumerate}

\subsection{Problem 6: Dielectric in a Parallel-Plate Capacitor}
A dielectric slab completely fills the space between parallel conducting plates held at fixed potential difference $V$.
\begin{enumerate}[label=\alph*)]
\item Determine the electric field in the dielectric.
\item Find the displacement field and free surface charge on the plates.
\item Compare the capacitance with the vacuum case.
\end{enumerate}

\section{Interfaces and Boundary Conditions}
\subsection{Problem 7: Interface Between Two Dielectrics}
Two linear dielectrics with permittivities $\varepsilon_1$ and $\varepsilon_2$ meet at a planar interface.
\begin{enumerate}[label=\alph*)]
\item State the boundary condition on the normal component of $\mathbf{D}$.
\item State the boundary condition on the tangential component of $\mathbf{E}$.
\item Explain how these conditions change if a free surface charge density is present.
\end{enumerate}

\subsection{Problem 8: Refraction of Field Lines}
Using the dielectric boundary conditions, derive the relation between the angles that the electric field makes with the normal on the two sides of a dielectric interface.

\section{Potential Formulation}
\subsection{Problem 9: Poisson Equation in Dielectrics}
Starting from
\[
\nabla\cdot(\varepsilon \nabla \Phi)=-\rho_f,
\]
derive the form of the electrostatic equation for
\begin{enumerate}[label=\alph*)]
\item constant $\varepsilon$,
\item piecewise constant $\varepsilon$,
\item no free charge.
\end{enumerate}

\subsection{Problem 10: Piecewise Constant Permittivity}
In one dimension, let
\[
\varepsilon(x)=
\begin{cases}
\varepsilon_1, & x<0,\\
\varepsilon_2, & x>0.
\end{cases}
\]
Assume no free charge except possibly at the interface. State the continuity or jump conditions satisfied by $\Phi$ and $\varepsilon\,d\Phi/dx$ at $x=0$.

\section{Dielectric Sphere in an External Field}
\subsection{Problem 11: Polarizable Sphere}
A dielectric sphere of radius $R$ and permittivity $\varepsilon$ is placed in an initially uniform external field $\mathbf{E}_0=E_0\hat{\mathbf{z}}$.
\begin{enumerate}[label=\alph*)]
\item Write the general form of the interior and exterior potentials.
\item Impose the boundary conditions at $r=R$.
\item State the resulting interior field and induced dipole moment.
\end{enumerate}

\subsection{Problem 12: Limiting Cases}
For the dielectric sphere in the previous problem, discuss the limits:
\begin{enumerate}[label=\alph*)]
\item $\varepsilon=\varepsilon_0$,
\item $\varepsilon\to \infty$,
\item $\varepsilon \gg \varepsilon_0$ but finite.
\end{enumerate}

\section{Energy and Interpretation}
\subsection{Problem 13: Electrostatic Energy in Matter}
Write the electrostatic energy density in a linear dielectric. Explain why the presence of dielectric material changes the energy stored at fixed free charge.

\subsection{Problem 14: Microscopic Versus Macroscopic Picture}
Explain the conceptual difference between microscopic charges and the macroscopic fields $\mathbf{P}$ and $\mathbf{D}$. Why are bound charges useful even though they are not independent sources?

\section{Challenge Problems}
\subsection{Problem 15: Point Charge in a Uniform Dielectric}
A point charge $q$ is embedded in a homogeneous linear dielectric of permittivity $\varepsilon$.
\begin{enumerate}[label=\alph*)]
\item Write the potential.
\item Compare it with the vacuum result.
\item Explain the physical meaning of the reduction by $\varepsilon$.
\end{enumerate}

\subsection{Problem 16: Common Pitfall}
A student says that if $\rho_f=0$ then $\mathbf{D}=0$ everywhere. Explain why this is false and give a counterexample.

\end{document}
